% !PRINTED 458
\weberpage{458}

If, conversely, such an equation holds, then by comparison with (4)
there follows, in the manner applied earlier, when $x, y$ are
integers:

\weq{7}{
\begin{aligned}
\partial\beta &= Ax, \\
c\alpha-a\gamma &= Cx, \\
c\beta+\partial\alpha-a\delta &= Bx, \\
c\beta-\partial\alpha-a\delta &= y,
\end{aligned}
}

from which:

\weq{8}{
\begin{aligned}
2\partial\alpha &= Bx-y, \\
2(c\beta-a\delta) &= Bx+y, \\
4n &= y^2-Dx^2.
\end{aligned}
}

If we assume $A$ relatively prime to $3n$, then $\partial = 1$,
$a = n$ must be taken; for by (7), under this assumption $x$ must be
divisible by $\partial$, and consequently also $c\alpha-a\gamma$ and
$c\beta-a\delta$, from which it follows that $a, c$ too would have to
be divisible by $\partial$, whereas $a, c, \partial$ have no common
divisor. Thus

\weq{9}{
\begin{aligned}
2\alpha &= Bx-y, & \beta &= Ax, \\
2n\gamma &= c(Bx-y)-2Cx, & 2n\delta &= 2cAx-Bx-y.
\end{aligned}
}

Besides the restrictions stated in \S\,114, 1), 2), the numbers
$x, y$ are here subject, on account of (3), to the further condition
that $y^2-Dx^2 \equiv -1\ (\mathrm{mod}\ 3)$. \gesperrt{This excludes the
possibility that $D$ is divisible by 3.} $x$ must be different from
zero, and we may assume it positive without restricting the
generality. (A simultaneous change of sign of $x$ and $y$ produces
only a change of sign of $\alpha, \beta, \gamma, \delta$, which is of
no consequence.)

If $D \equiv -1\ (\mathrm{mod}\ 3)$, then $x$ and $y$ must be
indivisible by 3; if $D \equiv 1\ (\mathrm{mod}\ 3)$, then $y$ is
divisible by 3 and $x$ is not divisible. For the rest, the numbers
$x, y$ may be assumed arbitrarily. Then $c$ is determined from the
two (mutually compatible) congruences:

\[
\begin{aligned}
cAx &\equiv \frac{Bx+y}{2} \\[4pt]
c\,\frac{Bx-y}{2} &\equiv Cx
\end{aligned}
\quad (\mathrm{mod}\ n)
\]

and, since $n$ is indivisible by 3, it can be subjected to the
condition

\[
c \equiv 0\ (\mathrm{mod}\ 3)
\]

as well.
